\documentclass[aps, amsmath,amssymb, reprint, twocolumn, superscriptaddress]{revtex4-1}
\usepackage[utf8]{inputenc}

\setcitestyle{super}

\usepackage{graphicx}
\usepackage{dcolumn}
\usepackage{bm}
\usepackage{amsmath}
\usepackage{epstopdf}
\usepackage{epsfig}
\usepackage[caption=false]{subfig}
\usepackage{fancyvrb}
\usepackage[usenames,dvipsnames]{color}
\usepackage[normalem]{ulem}
\usepackage{mathbbol}
\usepackage{tikz}
\usepackage{gensymb}
\usepackage{booktabs}
\usepackage{placeins}   % provides \FloatBarrier

\newcommand{\br}{\mathbf{r}}
\newcommand{\bk}{\mathbf{k}}
\newcommand{\dd}[2]{\frac{\partial #1}{\partial #2}}

\newcommand{\comment}[1]{\textcolor{Mahogany}{#1}}
\newcommand{\edit}[1]{\textcolor{OrangeRed}{#1}}
\newcommand{\delete}[1]{\textcolor{OrangeRed}{\sout{#1}}}
\newcommand{\BeginEdit}{\color{OrangeRed}}
\newcommand{\EndEdit}{\color{black}}

%%%%%%%%%%%%%%%%%%%%%%%%%%%%%%%%%%%%%%%%%%%%%%%%%%%%%%%%%%

\begin{document}

\title{Is Pure Water a Mixture: Perspectives from Machine-Learning Classification Methods}
\author{Yiqi Yao}
\affiliation{Department of Chemistry, Harvey Mudd College, 301 Platt Boulevard, Claremont, California 91711}
\author{Diya Sanghi}
\affiliation{Department of Chemistry, Harvey Mudd College, 301 Platt Boulevard, Claremont, California 91711}
\author{Akanksha D. Chokshi}
\affiliation{Yale-NUS College, Singapore 138527, Singapore}
\author{Minwoo Choi}
\affiliation{Yale-NUS College, Singapore 138527, Singapore}
\author{Haiqin Wang}
\affiliation{Department of Chemistry, Harvey Mudd College, 301 Platt Boulevard, Claremont, California 91711}
\author{Bilin Zhuang}
\email{bzhuang@hmc.edu}
\affiliation{Department of Chemistry, Harvey Mudd College, 301 Platt Boulevard, Claremont, California 91711}
\affiliation{Yale-NUS College, Singapore 138527, Singapore}
\date{\today}

\begin{abstract}
Whether liquid water is a structural continuum or a dynamic mixture of two distinct local environments 
has remained a matter of serious debate for over a century. By applying unsupervised machine-learning
clustering to molecular dynamics simulations of TIP4P/2005 and TIP5P water, we show that two
structurally distinct local populations---locally favored tetrahedral structures (LFTS) and disordered
normal-liquid structures (DNLS)---can be recovered without prior knowledge of the classification scheme.
A systematic comparison of clustering algorithms operating on a five-dimensional order-parameter feature
space ($q$, LSI, $S_k$, $\zeta$) reveals that hybrid density-denoising pipelines yield the clearest
two-state separation. Crucially, the cluster assignments are validated through an independent physical
observable: the per-cluster oxygen--oxygen partial structure factor $S(k)$, computed via the Debye
scattering equation directly from atomic coordinates. The LFTS cluster exhibits a first sharp diffraction
peak at $k_{T1}\approx 0.81$, while the DNLS cluster peaks at $k_{D1} \approx 1.05$, in quantitative
agreement with the theoretical predictions of the two-state model. Because the clustering and validation
share no common descriptor, this agreement constitutes model-independent confirmation that the two-state
signature in the structure factor reflects genuine structural heterogeneity in liquid water rather than an
artifact of the classification scheme.
\end{abstract}

\maketitle
% 1. Introduction

\section{Introduction}

The local structure of liquid water has been argued about for over a century. Two fundamentally different pictures have competed to explain water's thermodynamic and dynamic anomalies, including the density maximum at 4$^\circ$C, the rapid increase in isothermal compressibility upon supercooling, and the fragile-to-strong dynamic transition~\cite{gallo_water_2016}. Continuum models describe water through a broad unimodal distribution of local environments~\cite{narten_observed_1969, eisenberg_structure_2005}. Mixture models, dating back to R\"{o}ntgen's 1892 hypothesis of coexisting ``icebergs'' and a fluid ``sea''~\cite{rontgen-1892}, instead posit two or more distinct structural populations whose relative fractions shift with temperature and pressure. Both pictures can account for many of water's well-known anomalies, which is precisely why the debate has proven so difficult to settle. Despite a clear difference in
physical picture, unambiguous experimental evidence has remained elusive~\cite{narten_observed_1969,eisenberg_structure_2005,handle_supercooled_2017}.

Recent computational work by Shi and Tanaka has given the two-state picture considerably firmer footing~\cite{shi_direct_2020,russo_understanding_2014,shi_microscopic_2018}. In their framework, which we refer to hereafter as the two-state model, water is a dynamic mixture of locally favored tetrahedral structures (LFTS), stabilized by four hydrogen bonds with low density and high local symmetry, and disordered normal-liquid structures (DNLS), characterized by broken tetrahedral symmetry, higher density, and greater entropy. A key tool in their analysis is the structural descriptor $\zeta$, which quantifies translational order in the second coordination shell~\cite{russo_understanding_2014}. They demonstrated that the distribution $P(\zeta)$ is bimodal in several classical water models, and that the oxygen--oxygen partial structure factor $S(k)$ contains two overlapping peaks within the apparent first diffraction peak: a first sharp diffraction peak (FSDP) at $k_{T1} = kr_{\mathrm{OO}}/2\pi \approx 3/4$ arising from the tetrahedral LFTS geometry, and an ordinary liquid peak at $k_{D1} \approx 1$ characteristic of DNLS~\cite{shi_direct_2020,shi-2019}.

In this work, unsupervised machine-learning algorithms are applied to a four-dimensional order-parameter feature vector ($q$, LSI, $S_k$, $\zeta$) computed from molecular dynamics trajectories of TIP4P/2005 and TIP5P water. Four clustering strategies are benchmarked---K-Means, Gaussian Mixture Models (GMM), and two hybrid pipelines combining density-based denoising (DBSCAN, HDBSCAN) with GMM---and the resulting cluster assignments are validated through the per-cluster oxygen--oxygen partial structure factor $S(k)$, computed via the Debye scattering equation directly from atomic coordinates. The key point is that clustering is performed in order-parameter space while validation uses a reciprocal-space quantity derived from atomic coordinates.


\FloatBarrier
% 2. Method



\section{Results and Discussion}

\subsection{Overall Classification Workflow}

In the following section, we benchmark all algorithms on TIP4P/2005 model at T=253.15T = 253.15 $T = 253.15$~K ($-20\,^\circ$C), near the Schottky temperature $T_{s=1/2} \approx 237.8$~K where LFTS and DNLS populations coexist in appreciable proportions.

\subsection{Classification Algorithm Comparison}

\subsubsection{Baseline: K-Means and GMM}

K-Means ($k=2$) divides the $N = 20{,}480$ molecular configurations into Cluster~0 (11807 molecules, 57.7\%) and Cluster~1 (8673 molecules, 42.3\%). The per-cluster order-parameter distributions and pairwise projections (Fig.~\ref{fig:clustering_comparison}, panels a--b) show that the two clusters correspond to LFTS and DNLS: $\zeta$ produces well-separated bimodal peaks, and $q$ is higher on average in Cluster~0. The K-Means silhouette score is $0.2418$.

GMM yields qualitatively consistent populations (11661 and 8819 molecules) with a silhouette score of $0.2144$, slightly lower than K-Means, which is expected given its softer probabilistic boundaries. That both baseline methods recover the same two-state partition from the four-dimensional feature space suggests the signal is real and not an artifact of clustering method.\label{app:gmm}


% FIGURE 1 

\begin{figure*}[t]
  \centering
  \includegraphics[width=15cm]{Images/final_photos/fig1.png}
  \caption{%
   \textbf{ Comparison between the baseline K-Means and the hybrid clustering DBSCAN-GMM pipelines for classifying the TIP4P/2005 water at $T = 253.15$~K into locally-favored tetrahedral structures (LFTS) and disordered normal-liquid structures (DNLS).}
    \textbf{(a)} $q$--$\zeta$ plot for water molecules in the two states classified by the baseline K-Means pipeline, where $q$ and $\zeta$ are the tetrahedral order parameter and the local transitional order parameter, respectively. The silhouette score for this pipeline is $0.2418$.
    \textbf{(b)} Probability density functions $f(q)$, $f(\text{LSI})$, $f(S_k)$, and $f(\zeta)$ for order parameters $q$, $\text{LSI}$, $S_k$, and $\zeta$ that make up the dataset for water from the K-mean clustering. Here, LSI is the local structure index and $S_k$ is the transitional tetrahedral parameter. 
    \textbf{(c)} $q$--$\zeta$ plot for water molecules in the two states classified  by the hybrid clustering DBSCAN-GMM pipeline. 
    In addition to the two states, the pipeline flags noise (plotted in grey) indicating molecules that are in transition between two states. The silhouette score for this pipeline is $0.2841$.
    \textbf{(d)} Probability density functions $f(q)$, $f(\text{LSI})$, $f(S_k)$, and $f(\zeta)$ from the hybrid DBSCAN-GMM pipeline.
  }
  \label{fig:clustering_comparison}
\end{figure*}

\subsubsection{Hybrid DBSCAN--GMM Pipeline}
In the four-dimensional order-parameter space, LFTS and DNLS molecules each cluster into dense regions with a sparsely populated gap between them. DBSCAN treats molecules in this transition zone as noise, leaving a cleaner subset for GMM to work with. After the parameter sweep, the DBSCAN--GMM pipeline discards roughly $22\%$ of molecules as noise and reaches a silhouette score of $0.2841$---noticeably better than either baseline. The tighter cluster boundaries are visible in Fig.~\ref{fig:clustering_comparison}, panels c--d. We also tested HDBSCAN--GMM, which sidesteps the need to manually tune $\epsilon$ and removes a slightly smaller noise fraction ($15.2\%$), yielding a silhouette score of $0.2572$; full results in Appendix~\ref{app:hdbscan}.


\FloatBarrier
\subsection{Structure Factor Validation}
\label{sec:structure_factor_results}
We present the following structure factor validation steps with DBSCAN--GMM pipeline, given its highest silhouette score among the methods tested. As shown in Fig.~\ref{fig:sk_per_cluster}, the two clusters produce noticeably different $S(k)$ curves: the LFTS cluster peaks at $k_{T1} \approx 0.81$ while the DNLS cluster peaks at $k_{D1} \approx 1.05$, both highly agree with the predictions $k_{T1} \approx
3/4$ and $k_{D1} \approx 1$ of the two-state model~\cite{shi_direct_2020}.

% FIGURE 2 
\begin{figure*}[t]
  \centering
  \includegraphics[width=15cm]{Images/final_photos/fig2.png}
  \caption{%
    \textbf{Per-cluster oxygen--oxygen partial structure factors $S(k)$ for TIP4P/2005
    water at $T = 253.15$~K}, computed via the Debye scattering Eq.~\eqref{eq:debye} from the DBSCAN--GMM classification. The LFTS state displays its primary peak at
    $k_{T1} \approx 0.81\;(k\,r_{\mathrm{OO}}/2\pi)$, consistent with the first sharp diffraction peak (FSDP) of tetrahedral materials, while this peak is absent for the DNLS state, indicating significant tetrahedral feature in the LFTS state but not in the DNLS state.
  }
  \label{fig:sk_per_cluster}
\end{figure*}

The total structure factor shows a single broad peak sitting between $k_{T1}$ and $k_{D1}$, which indicates that the apparent first diffraction peak of liquid water is actually a superposition of two overlapping contributions from structurally distinct subpopulations. The $\zeta$-resolved structure factor $S(k,\zeta)$ (Fig.~\ref{fig:sk_zeta_surfaces}) adds further detail: two distinct ridges appear in separate $\zeta$ domains, with the FSDP at $k_{T1}$ concentrated in the high-$\zeta$ (LFTS) region and the ordinary peak at $k_{D1}$ dominating the low-$\zeta$ (DNLS) region, consistent with the surfaces reported by Shi and Tanaka (cf.\ Fig.~2D--E)~\cite{shi_direct_2020}.

% FIGURE 3
\begin{figure*}[t]
  \centering
  \includegraphics[width=19cm]{Images/final_photos/fig3.png}
  \caption{\textbf{$\zeta$-dependence in the structure factor for the LFTS and DNLS states for TIP4P/2005 water model
    at $T = 253.15$~K.}
    \textbf{(a,b)} Three-dimensional $S(k,\zeta)$ surfaces for the (a) the LFTS and (b) the DNLS states, respectively. 
    The blue circle highlight the first sharp diffraction peak (FSDP) at $k_{T1} \approx 3/4$ for tetrahedral materials
    concentrated in the high-$\zeta$ domain. The red circle highlight the
    ordinary liquid peak at $k_{D1} \approx 1$ concentrated in the low-$\zeta$ domain.
    \textbf{(c,d)} Corresponding two-dimensional contour maps.
    Vertical dashed lines mark $k_{T1}$ (blue, FSDP) and $k_{D1}$ (red, ordinary peak).
    Each cluster's $S(k)$ maximum corresponds to the other's minimum, confirming that the two structural populations are physically distinct.
  }
  \label{fig:sk_zeta_surfaces}
\end{figure*}

\FloatBarrier
\subsection{Robustness and Generality}

\subsubsection{Temperature Dependence}
We apply the DBSCAN--GMM pipeline to TIP4P/2005 water across seven temperatures from $-30\,^\circ$C to $+30\,^\circ$C (Fig.~\ref{fig:robustness}, panel a). The LFTS cluster responds to cooling with the FSDP at $k_{T1}$ strengthens monotonically as the LFTS fraction $s$ grows toward the low-density amorphous limit~\cite{shi-2019}. The DNLS cluster shows much weaker temperature dependence, with its peak near $k_{D1}$ remaining essentially stable throughout. This behavior is consistent with the two-state model and holds across the full supercooled range we examined.

\subsubsection{Model Dependence}
To check whether the two-state signatures depend on the choice of water model, we run the pipeline on three different water models at $T = -20\,^{\circ}$C. In the LFTS cluster (Fig.~\ref{fig:sk_temperature_model}c), all three produce a clear FSDP, pointing to tetrahedral ordering as a robust feature of the LFTS population. The peak intensities do differ, where TIP5P is the strongest, followed by TIP4P/2005, with SWM4-NDP noticeably weaker. The DNLS cluster (Fig.~\ref{fig:sk_temperature_model}d) is more uniform across models, with the ordinary liquid peak at $k_{D1} \approx 1.05$ staying stable throughout; SWM4-NDP shows the strongest DNLS signal of the three.



% FIGURE 4 — robustness

\begin{figure*}[t]
  \centering
  \includegraphics[width=15cm]{Images/final_photos/fig4.png}
  \caption{%
    \textbf{Temperature dependence of the per-cluster oxygen--oxygen partial
    structure factor $S(k)$}  for (a) the LFTS state and (b) the DNLS state. The vertical dashed lines marked the positions of the first sharp diffraction peak for tetrahedral material, $k_{T1}$, and the ordinary liquid peak, $k_{D1}$. 
    %\textbf{(c-d)} LFTS and DNLS cluster at $T = -20\,^{\circ}$C for three water models:
    %SWM4-NDP, TIP4P/2005, and TIP5P.
    %In all panels, vertical dashed lines mark the predicted peak positions of
    %the two structural populations: $k_{T1}$ (green, FSDP, characteristic of
    %LFTS tetrahedral ordering) and $k_{D1}$ (red, ordinary liquid peak,
    %characteristic of DNLS).
    %Error bars represent the standard error across 20 independent MD runs.
  }
  [expand on y axis, make the figure in y direction better]
  \label{fig:sk_temperature_model}
\end{figure*}



\FloatBarrier
%%%%%%%%%%%%
% 4. Conclusion
%%%%%%%%%%%%

\section{Conclusion}
\label{sec:conclusion}
Our results support a two-state picture of liquid water. Without any prior structural knowledge built in, our ML clustering recovers LFTS and DNLS as two distinct populations, and the DBSCAN--GMM pipeline does this most cleanly (silhouette score $0.2841$) by discarding the ambiguous transition-zone molecules before classification. More importantly, when we validate the cluster assignments through the oxygen--oxygen partial structure factor, the LFTS cluster peaks at $k_{T1} \approx 0.81$ and DNLS at $k_{D1} \approx 1.05$, matching the two-state model predictions. The convergence of ML clustering in order-parameter space and structure factor validation in reciprocal space provides a  compelling case for real structural heterogeneity in supercooled water, confirming the two-state model. 



\newpage
\input{Appendix}

\FloatBarrier
\section{Method}

\subsection{Molecular Dynamics Simulations}
To provide microscopic configurations of water molecules for our analysis, we carry out classical molecular dynamics simulations for systems of $N = 1024$ water molecules in the canonical (NVT) ensemble using the OpenMM simulation package~\cite{eastman_openmm_2015}  with the rigid, non-polarizable TIP4P/2005~\cite{abascal_general_2005} and TIP5P~\cite{mahoney_five-site_2000} models, which have demonstrated strong tetrahedral ordering signatures.\cite{shi_direct_2020} Simulations are conducted at $5\,^\circ$C increments from  $-30\,^\circ$C to $30\,^\circ$C. The density of water molecules in the system is set to be that of the experimentally observed densities at the corresponding temperature. The Schottky temperatures---at which the two populations are approximately equal---are $T_{s=1/2} \approx 237.8$~K for TIP4P/2005 and $T_{s=1/2} \approx 255.5$~K for TIP5P~\cite{shi_direct_2020}. For $T=-30\,^\circ$C to $30\,^\circ$C, we perform 20 independent simulations for each model using the Langevin integrator at 1 fs timestep with a friction coefficient 20 ps$^{-1}$. After equilibrating the system for 10 ns, the system is sampled every 10 ps for 50 ns. This yields $20480$ molecular configurations for each temperature and model. 

\subsection{Structural Order Parameters}
\label{sec:structural_order_parameter}
We compute the following four structural order parameters for each molecule to provide data that describe the local water arrangements: (i) orientational order parameter, $q$; (ii) local structure index, LSI; (iii) translational tetrahedral parameter, $S_k$; and (iv) local translational order parameter, $\zeta$. For every trajectory frame sampled from the MD simulation, we can compute the set of order parameters for each water molecule. For simplicity, the position of each water molecule is represented by its oxygen atom's position. The intermolecular distance between pairs of molecules are computed using MDTraj Package~\cite{mcgibbon_mdtraj_2015} with considerations of the periodic boundary conditions. 

The orientational order parameter, $q$, measures the extent to which the angular positions of four nearest neighbors of a water molecule form a tetrahedral arrangement. It is given by~\cite{chau_new_1998,errington_relationship_2001}
\begin{equation}
q = 1 - \frac{3}{8}\sum_{(j,k)}\left(\cos\psi_{jk}+\frac{1}{3}\right)^{2},
\end{equation}
where the sum runs through all pairs $(j,k)$ among the four nearest neighbors of a molecule. $\psi_{jk}$ is the angle extended by neighbors $j$ and $k$ from 

 is  between neighbors $j$ and $k$. The values range from $-3$ to $1$,with $q=1$ representing a perfect tetrahedron.

\item \textbf{Local structure index, LSI.}
LSI quantifies the gap between the first and second hydration shells surrounding each molecule~\cite{shiratani_growth_1996}:
\begin{equation}
\text{LSI} = \frac{1}{n}\sum_{i=1}^{n}\!\left(\Delta_i - \bar{\Delta}\right)^{2},
\label{eq:LSI}
\end{equation}
where $\Delta_i = r_{i+1} - r_i$, represents the distance gap between consecutively ordered oxygen neighbors within a cutoff of 3.7~\AA, and $\bar{\Delta}$ is the mean gap. A large LSI indicates well-separated hydration shells.

\item \textbf{Translational tetrahedral parameter, $S_k$.}
$S_k$ measures radial uniformity of the four nearest-neighbor distances~\cite{duboue-dijon_characterization_2015}:
\begin{equation}
S_k = 1 - \frac{1}{3}\sum_{k=1}^{4}\frac{(r_k - \bar{r})^{2}}{4\bar{r}^{2}},
\end{equation}
where $r_k$ is the distance to the $k\text{th}$ nearest neighbor and $\bar{r}$ is their mean. $S_k = 1$ for a perfect tetrahedron.

\item \textbf{Local translational order parameter, $\zeta$.}
Russo and Tanaka~\cite{russo_understanding_2014} introduced $\zeta$ to characterize translational order in the second coordination shell:
\begin{equation}
\zeta = r_{\text{nearest non-HB}} - r_{\text{farthest HB}},
\label{eq:zeta}
\end{equation}
where $r_{\text{farthest HB}}$ is the distance to the most distant hydrogen-bonded neighbor and $r_{\text{nearest non-HB}}$ is the distance to the closest non-hydrogen-bonded neighbor. Large positive $\zeta$ indicates clear shell separation (LFTS); $\zeta \approx 0$ or negative indicates shell interpenetration (DNLS).

\

We computed all four parameters per molecule per frame and stored as a feature matrix of dimensions $N_{\text{samples}} \times 4$, where $N_{\text{samples}} = N_{\text{frames}} \times N_{\text{molecules}} \times N_{\text{runs}}$. We standardize the four order parameters by scaling the values to the unit interval $[0,1]$ prior to clustering. 

\subsection{Clustering Methods}
\label{sec:clustering_algorithm}
%Six unsupervised clustering methods are considered in this work, and they fall into three general categories as described in Ref.~\citenum{zheng-2018}
%applied six unsupervised methods, falling into three categories~\cite{zheng-2018}.

\subsubsection{Baseline: K-Means and GMM}
We apply K-Means~\cite{ikotun_k-means_2023}, which partitions the standardized feature space into two disjoint clusters ($k=2$) by iteratively assigning molecules to the nearest  centroid. We initialize centroids via the K-Means++ algorithm [CITE] available in the sckit-learn package\cite{scikit-learn}, with $n_{\text{init}} = 20$ restarts and a 300 iteration maximum. Building on findings that two-state frameworks yield nearly Gaussian  order-parameter distributions~\cite{shi_direct_2020}, we also fit a GMM as a more robust  probabilistic alternative~\cite{deprez-2024}. We fit GMMs with $k=2$, and additionally evaluate  models with $k=2,3,4$ components using BIC and AIC to test for the presence of additional  structural states.

\subsubsection{Hybrid Density--GMM Pipelines}
DBSCAN~\cite{ester_density-based_1996} identifies clusters as contiguous high-density regions and marks low-density boundary molecules as noise. In hybrid pipelines, we applied DBSCAN as a denoising step since molecules in the low-density transition zone between the two structural states are identified and removed. Then,  we fit GMM to the denoised subset. This two-stage approach combines the noise-identification capability of density-based algorithms with GMM, yielding substantially higher silhouette scores. The DBSCAN hyperparameter $\varepsilon$ (neighborhood radius in units of the standardized feature space) and MinPts (minimum cluster size) are selected via a systematic grid search detailed in Appendix~\ref{app:parameter_summary}.


\subsection{Structural Validation via Structure Factor Calculation}
\label{sec:structure_factor}
The two-state model demonstrated that the $\zeta$-resolved structure factor $S(k,\zeta)$ reveals two distinct peaks, thereby establishing a powerful connection between microscopic structural descriptors and experimentally accessible scattering signatures. We build on this foundation by asking whether the same reciprocal-space signatures emerge when we derive cluster assignments through a multi-dimensional ML pipeline. Specifically, we map cluster labels back onto the MD trajectories and compute the oxygen--oxygen partial structure factor separately for each cluster subpopulation using the Debye scattering equation~\cite{debye_zerstreuung_1915}
\begin{equation}
S(k) = 1 + \frac{1}{N}\sum_{i=1}^{N}\sum_{\substack{j=1 \\ j\neq i}}^{N}\frac{\sin(k r_{ij})}{k r_{ij}}\,W(r_{ij})
\label{eq:debye}
\end{equation}
where $W(r_{ij}) = \sin(\pi r_{ij}/r_c)/(\pi r_{ij}/r_c)$ is a sinc window function that suppresses Gibbs ringing artifacts from the hard distance cutoff $r_c = 1.5$~nm. Because our clustering draws on a multi-dimensional feature space without privileging any single order parameter, successfully recovering the characteristic $k_{T1}$ and $k_{D1}$ peaks would provide an complementary confirmation of the two-state model.

\FloatBarrier
% 3. Results and Discussion


\section*{Bibliography}
\bibliographystyle{apsrev4-2}
\bibliography{bib}

\end{document}
